%
% =================================================================
%  Template LaTeX Tesis Program Magister
%  Fakultas Matematika dan Ilmu Pengetahuan Alam (FMIPA)
%  Universitas Syiah Kuala (USK)
% =================================================================
%
%  Penulis       : Muhammad Ichsan
%  NPM           : 260820701100007
%  Hak cipta     : (c) 2026 Muhammad Ichsan
%                  Blok ini hanya penanda kepemilikan dan tidak
%                  muncul pada halaman hasil cetak.
%
%  Dasar         : Mengikuti Panduan Penulisan Tugas Akhir dan Tesis
%                  FMIPA Universitas Syiah Kuala tahun 2024.
%                  Seluruh struktur bab, halaman depan, dan daftar isi
%                  disusun sesuai panduan tersebut.
%
%  Terakhir diperbarui : 8 Oktober 2026
% =================================================================
%
% File ini adalah ISI UTAMA naskah. Jangan dikompilasi langsung;
% gunakan salah satu file pemanggil berikut yang menentukan tahap:
%
%   tesis.tex              \def\Tahap{finals}   -> tesis.pdf
%   seminar-proposal.tex   \def\Tahap{sempro}   -> seminar-proposal.pdf
%   seminar-hasil.tex      \def\Tahap{semhas}   -> seminar-hasil.pdf
%
% Cara memakai:
%   1. Isi data pribadi di include/identitas.tex
%   2. Tulis naskah pada file bab1.tex sampai bab5.tex serta lampiran.tex
%   3. Jalankan ./build.sh [sempro|semhas|sidang|all]  (Linux/macOS)
%      atau .\build.bat [sempro|semhas|sidang|all]     (Windows)
%   4. Kompilasi dilakukan 3 kali agar daftar isi dan nomor halaman stabil
%

% Karena class file berada di folder lib, tambahkan path lib/ ke path pencarian
\makeatletter
\def\input@path{{lib/}}
\makeatother
\documentclass[dvipsnames]{jiftesistheme}

\tolerance=1
\emergencystretch=\maxdimen
\hyphenpenalty=10000
\hbadness=10000

% File hype.indonesia.tex berada di folder language
\makeatletter
\def\input@path{{language/}}
\makeatother
% Daftar pemenggalan suku kata dan istilah dalam LaTeX
\input{hype.indonesia}

% Untuk prefiks pada daftar gambar dan tabel
\usepackage[titles]{tocloft}

\usepackage{etoolbox}
\makeatletter
% \patchcmd{<cmd>}{<search>}{<replace>}{<succes>}{<failure>}
\patchcmd{\@chapter}{\addtocontents{lof}{\protect\addvspace{10\p@}}}{}{}{}% LoF
\patchcmd{\@chapter}{\addtocontents{lot}{\protect\addvspace{10\p@}}}{}{}{}% LoT
\makeatother

\usepackage[justification=centering]{caption}

% Pengaturan huruf tebal pada daftar isi
\newcommand*{\enableboldchapterintoc}{%
  \addtocontents{toc}{\string\renewcommand{\protect\cftchapfont}{\protect\normalfont\protect\bfseries}}
  \addtocontents{toc}{\string\renewcommand{\protect\cftchappagefont}{\protect\normalfont\protect}}
  \addtocontents{toc}{\protect\setlength{\cftbeforechapskip}{12pt}}
}
\newcommand*{\disableboldchapterintoc}{%
  \addtocontents{toc}{\string\renewcommand{\protect\cftchappagefont}{\protect\normalfont}}
  \addtocontents{toc}{\string\renewcommand{\protect\cftchapfont}{\protect\normalfont}}
  \addtocontents{toc}{\protect\setlength{\cftbeforechapskip}{0pt}}
}

\renewcommand{\cftdotsep}{0.5}
\renewcommand{\cftchapleader}{\cftdotfill{\cftdotsep}}

\renewcommand\cftfigpresnum{Gambar\  }
\renewcommand\cfttabpresnum{Tabel\   }

\newcommand{\listappendicesname}{DAFTAR LAMPIRAN}
\newlistof{appendices}{apc}{\listappendicesname}
\newcommand{\appendices}[1]{\addcontentsline{apc}{appendices}{#1}}
\newcommand{\newappendix}[1]{\section*{#1}\appendices{#1}}

% Untuk hyperlink dan daftar isi
\usepackage[hidelinks]{hyperref}
\renewcommand\UrlFont{\rmfamily\itshape}
\newlength{\mylenf}
\settowidth{\mylenf}{\cftfigpresnum}
\setlength{\cftfignumwidth}{\dimexpr\mylenf+2em}
\setlength{\cfttabnumwidth}{\dimexpr\mylenf+2em}

% Agar ada tulisan BAB pada daftar isi
\renewcommand\cftchappresnum{BAB }
  \cftsetindents{chapter}{0em}{4.5em} % indentasi bab
  \cftsetindents{section}{4.5em}{2em}
  \cftsetindents{subsection}{6.5em}{3em}

% Agar di daftar isi setiap nomor subbab diakhiri titik
\renewcommand{\cftsecaftersnum}{.}
\renewcommand{\cftsubsecaftersnum}{.}

\addtocontents{toc}{~\hfill \textit{Halaman}\par} % Menambahkan kata "Halaman" di kanan daftar isi

% Membuat juga kata "Halaman" di daftar gambar, tabel, dan lampiran
\addtocontents{lof}{~\hfill \textit{Halaman}\par}
\addtocontents{lot}{~\hfill \textit{Halaman}\par}
\addtocontents{apc}{~\hfill \textit{Halaman}\par}

\usepackage{titlesec}
\titlelabel{\thetitle.\quad}

% Agar caption tabel dan gambar diakhiri titik dan label tebal
\usepackage[labelsep=period]{caption}
\usepackage[labelfont=bf]{caption}
\usepackage{caption}
\usepackage{subcaption}

% Warna LaTeX
\usepackage{color}

% Tabel panjang dan halaman landscape
\usepackage{longtable}
\usepackage{pdflscape}
\usepackage{lscape}

% Cuplikan kode
\usepackage{listings, lstautogobble}

\usepackage{adjustbox}
\usepackage{fancybox, graphicx}
\usepackage{url}
\usepackage{microtype}
\usepackage{siunitx}
\usepackage{xcolor}
\usepackage{multirow}
\usepackage[normalem]{ulem}
\useunder{\uline}{\ul}{}

\usepackage{array}
\newcolumntype{P}[1]{>{\centering\arraybackslash}p{#1}}
\newcolumntype{M}[1]{>{\centering\arraybackslash}m{#1}}

\makeatletter
\def\input@path{{include/}}
\makeatother

% Awal dokumen
\usepackage{fancyhdr}
\usepackage{rotating}

% Prefiks pada Daftar Program
\makeatletter
\begingroup\let\newcounter\@gobble\let\setcounter\@gobbletwo
\globaldefs\@ne \let\c@loldepth\@ne
\newlistof{listings}{lol}{\lstlistlistingname}
\endgroup
\let\l@lstlisting\l@listings
\AtBeginDocument{\addtocontents{lol}{\protect\addvspace{10\p@}}}
\makeatother
\renewcommand{\lstlistoflistings}{\listoflistings}
\renewcommand\cftlistingspresnum{Program~}
\cftsetindents{listings}{1.5em}{7em}

% Mengatur indentasi daftar pustaka
\setlength{\bibhang}{30pt}

\usepackage{amsmath}
\usepackage{pdfpages}
\usepackage[table,xcdraw]{xcolor}
\usepackage{colortbl}

\begin{document}

%-----------------------------------------------------------------
% DATA IDENTITAS (isi di include/identitas.tex)
%-----------------------------------------------------------------
\input{identitas}

%-----------------------------------------------------------------
% HALAMAN DEPAN
% Urutan halaman depan mengikuti panduan penulisan tesis:
%   Halaman Judul, Halaman Pengesahan, [Pernyataan Bebas Plagiasi],
%   Kata Pengantar, Abstrak, Abstract, Daftar Isi, Daftar Tabel,
%   Daftar Gambar, Daftar Lampiran
%-----------------------------------------------------------------
\fancyhf{}
\fancyfoot[C]{\thepage}

\cover

\approvalpage

%-----------------------------------------------------------------
% Naskah finals juga memuat pernyataan bebas plagiasi
%-----------------------------------------------------------------
\iffinals
  \include{bebas-plagiasi}
\fi

%-----------------------------------------------------------------
% Surat pernyataan belum dipakai pada tesis ini.
% Halaman ini sengaja dinonaktifkan (tidak dihapus). Untuk
% mengaktifkan kembali, hapus tanda komentar pada baris di bawah.
% -----------------------------------------------------------------
%\include{surat-pernyataan}

\include{kata-pengantar}

\titleformat{\section}{\normalfont\bfseries\uppercase}{\thesection}{1.7em}{} % Section kapital, daftar isi tetap non kapital
\titleformat{\subsection}{\normalfont\bfseries}{\thesubsection}{0.9em}{} % Subsection jarak 0.9em

\include{abstrak-indonesia}

\include{abstrak-english}

%-----------------------------------------------------------------
% DAFTAR ISI, DAFTAR TABEL, DAFTAR GAMBAR, DAFTAR LAMPIRAN
%-----------------------------------------------------------------
\addcontentsline{toc}{chapter}{Daftar Isi}
\begin{singlespace}
	\tableofcontents
\end{singlespace}
\listoftables
\addcontentsline{toc}{chapter}{Daftar Tabel}
\listoffigures
\addcontentsline{toc}{chapter}{Daftar Gambar}

% Daftar Program (aktifkan bila tesis memuat cuplikan kode)
%\lstlistoflistings
%\addcontentsline{toc}{chapter}{Daftar Program}

% Daftar Lampiran hanya ada pada naskah finals
\iffinals
\listofappendices
\addcontentsline{toc}{chapter}{Daftar Lampiran}
\fi

\enableboldchapterintoc

%-----------------------------------------------------------------
% AWAL MASUKAN BAB
%-----------------------------------------------------------------
\begin{onehalfspace}

\fancyhf{}
\fancyfoot[R]{\thepage}

% PENTING: \clearpage dipakai agar halaman Daftar Lampiran (atau
% Daftar Gambar pada tahap Proposal) sudah selesai sebelum penomoran
% arabic dimulai. Tanpa ini halaman tersebut ikut bernomor arabic.
\clearpage
\pagenumbering{arabic}

\captionsetup[figure]{labelfont={normalfont}, textfont={normalfont}} % Caption gambar tidak tebal
\captionsetup[table]{labelfont={normalfont}, textfont={normalfont}} % Caption tabel tidak tebal

% Caption untuk cuplikan kode
\renewcommand{\thelstlisting}{\arabic{chapter}.\arabic{lstlisting}}
\renewcommand*\lstlistingname{Program}

\include{bab1}

\include{bab2}

\include{bab3}

% Seminar proposal belum memiliki Bab IV Hasil dan Pembahasan, sehingga
% Bab IV dan Bab V baru ikut sejak tahap seminar hasil.
\ifnotsempro
  \include{bab4}

  \include{bab5}
\fi

%-----------------------------------------------------------------
% DAFTAR PUSTAKA
%-----------------------------------------------------------------
\fancypagestyle{daftarpustaka}{
    \fancyhf{}
    \fancyfoot[R]{\thepage}
    \renewcommand{\headrulewidth}{0pt}
    \renewcommand{\footrulewidth}{0pt}
}

\addcontentsline{toc}{chapter}{DAFTAR PUSTAKA}
\begin{spacing}{1}
\pagestyle{daftarpustaka}
\bibliography{daftar-pustaka}
\end{spacing}

%-----------------------------------------------------------------
% LAMPIRAN (khusus naskah finals)
%-----------------------------------------------------------------
\iffinals
\include{lampiran}
\fi

\end{onehalfspace}

\end{document}